\documentclass{article}
\usepackage[T1]{fontenc}
\usepackage{lmodern}
\usepackage{graphicx}
\usepackage{amsmath,amssymb}
\usepackage[a4paper,margin=2cm]{geometry}
\usepackage[absolute,overlay]{textpos}
\setlength{\TPHorizModule}{1mm}
\setlength{\TPVertModule}{1mm}
\usepackage[indonesian]{babel}
\usepackage{hyperref}
\setlength{\emergencystretch}{2em}
% unit: Halaman 3

\begin{document}

% === Halaman 3 (llm) ===

\begin{itemize}
  \item ECC adalah kriptografi kunci-publik yang relatif lebih baru usianya.
  \item Dikembangkan oleh Neal Koblitz dan Victor S. Miller tahun 1985.
  \item Klaim: Panjang kunci ECC lebih pendek daripada kunci RSA, namun memiliki tingkat keamanan yang sama dengan RSA.
  \item Contoh: kunci ECC sepanjang 160-bit menyediakan tingkat keamanan yang sama dengan 1024-bit kunci RSA.
  \item Keuntungan: dengan panjang kunci yang lebih pendek, membutuhkan memori dan komputasi yang lebih sedikit.
  \item Cocok untuk piranti nirkabel, dimana prosesor, memori, umur batere terbatas.
\end{itemize}

\end{document}
